\documentclass{amsart}
\usepackage{amsmath,amssymb,amsthm,hyperref}
\newtheorem{theorem}{Theorem}
\newtheorem{proposition}{Proposition}
\newtheorem{conjecture}{Conjecture}
\newtheorem{remark}{Remark}
\begin{document}

\title{On sixth powers as a sum of five sixth powers}
\maketitle

\section*{Statement}
We are asked whether there exist natural numbers
$x_1,x_2,x_3,x_4,x_5,z\in\mathbb{N}$ with
\begin{equation}\label{eq:main}
x_1^6+x_2^6+x_3^6+x_4^6+x_5^6=z^6 .
\end{equation}

\section*{Honest status of the problem}

This is, to the best of current mathematical knowledge, an
\textbf{open problem}. No solution of \eqref{eq:main} in positive
integers is known, and no proof that \eqref{eq:main} has no solution
is known. In what follows we record precisely what is established, so
that no unproven claim is presented as a theorem.

\subsection*{Context: the Lander--Parkin--Selfridge framework}
For a fixed exponent $k$, write $(k.m.n)$ for a Diophantine equation
$$\sum_{i=1}^{m}a_i^{k}=\sum_{j=1}^{n}b_j^{k}$$
in positive integers. Equation \eqref{eq:main} is an instance of
$(6.5.1)$, i.e.\ $m=5$, $n=1$, $k=6$, so $m+n=6=k$.

\begin{conjecture}[Lander--Parkin--Selfridge, 1967]
If $\sum_{i=1}^{m}a_i^{k}=\sum_{j=1}^{n}b_j^{k}$ holds in positive
integers with all bases nonzero, then $m+n\ge k$.
\end{conjecture}

For \eqref{eq:main} we have $m+n=k=6$, exactly on the conjectural
boundary. The Lander--Parkin--Selfridge conjecture does \emph{not}
forbid \eqref{eq:main}; it only forbids representations with $m+n<6$,
e.g.\ a sixth power as a sum of four or fewer positive sixth powers.
Thus even a proof of their conjecture would leave \eqref{eq:main}
unresolved. Conversely, the conjecture itself is wide open.

\subsection*{What is known for sixth powers}
The smallest \emph{known} representation of a sixth power as a sum of
positive sixth powers uses \emph{seven} terms. Explicitly, one has the
identity
\begin{equation}\label{eq:seven}
74^6+234^6+402^6+474^6+702^6+894^6+1077^6=1141^6 ,
\end{equation}
which is verified by direct computation (both sides equal
$2206550475483180841$). No representation of a sixth power as a sum of
six or fewer positive sixth powers is currently known. In particular,
the five-term case \eqref{eq:main} has no known solution.

\subsection*{Bounded computational evidence}
A direct search establishes that \eqref{eq:main} has no solution with
small bases.

\begin{proposition}\label{prop:search}
Equation \eqref{eq:main} has no solution in positive integers with
$\max(x_1,\dots,x_5)\le 120$.
\end{proposition}

\begin{proof}[Proof (finite verification)]
Order the variables $x_1\le x_2\le\cdots\le x_5$. For every value of
$x_5\in\{1,\dots,120\}$ and every value of the partial sum
$s_1=x_1^6+x_2^6$ realizable with $x_1\le x_2\le 120$, one checks each
sixth power $z^6\ge x_5^6+2s_1$ and tests whether the residual
$$s_2:=z^6-x_5^6-s_1$$
is itself a sum of two sixth powers $x_3^6+x_4^6$ with
$s_2\ge s_1$ (so that the ordering constraint can be met). Since the
set of two-term sums $\{a^6+b^6: 1\le a\le b\le 120\}$ is finite and
precomputable, this is a finite exhaustive check over all admissible
$(x_5,s_1,z)$. The check returns no solution. (The same procedure was
also run, partially, for bases up to $250$ without finding any
solution.)
\end{proof}

\section*{Conclusion}

We do \textbf{not} have a proof either way for the general question
posed. The honest answer is:
\begin{itemize}
\item No tuple $(x_1,\dots,x_5,z)$ satisfying \eqref{eq:main} is known,
      and an exhaustive search over $\max x_i\le 120$ (Proposition
      \ref{prop:search}) finds none.
\item No proof of non-existence is known; the relevant
      Lander--Parkin--Selfridge conjecture, which in any case does not
      rule out the boundary case $(6.5.1)$, is itself unproven.
\item The current world record for representing a sixth power as a sum
      of positive sixth powers uses seven terms, as in \eqref{eq:seven};
      five terms remain out of reach in both directions.
\end{itemize}

Therefore the problem as stated is open, and this document records the
verified facts rather than asserting an unproven theorem.

\end{document}
